\documentclass[12pt,a4paper,oneside]{article}
\usepackage[brazil]{babel}
\usepackage[utf8]{inputenc}
\usepackage{olymp}
\usepackage{color}
\usepackage[dvipsnames]{xcolor}
\usepackage{listings}

\setcounter{problem}{3}

\begin{document}

\begin{problem}{Mediana}{mediana}{2}

Joãozinho acabou de aprender que mediana não é a mesma coisa que média. Mediana é o valor que separa a metade superior da metade inferior de um conjunto de dados. Em outras palavras, é o elemento ``do meio''. Por exemplo, a mediana dos elementos $\{1, 1, 2, 3, 5, 7, 9\}$ é o valor $3$. Diferente da média, que
nesse caso é $4$. Quando o conjunto de dados tem tamanho par, a mediana pode ser tomada como a média dos elementos centrais. Por exemplo, a mediana dos elementos $\{1, 1, 2, 3, 5, 7\}$ é dada por $(2+3)/2 = 2,5$.

A dificuldade de Joãozinho em encontrar a mediana acontece quando os elementos não estão em ordem. Faça um programa para ajudá-lo.

%\Notes
%
%Observações, se tiver.

\InputFile

Cada entrada contém apenas um caso de teste. A primeira linha contém um número $N$, indicando o número de elementos do conjunto. A segunda linha contém $N$ números inteiros, que são os elementos do conjunto.  Restrições: $1 \leq N \leq 1000$, $0 \leq$ cada elemento $\leq 100000$.

\OutputFile

Seu programa deve gerar 2 linhas na saída, a primeira contendo a média e a segunda contendo a mediana dos $N$ elementos do conjunto. Estes valores devem ser escritos com 1 casa decimal e devem seguir o formato mostrado nos exemplos.

%  \Notes
% Você pode escrever os resultados à medida que lê a entrada, não precisa ler a entrada toda antes.

\Example

\begin{example}
\exmp{
7
1 1 2 3 5 7 9
}{
Media: 4.0
Mediana: 3.0
}%
\end{example}

\begin{example}
\exmp{
7
1 3 5 9 7 1 2
}{
Media: 4.0
Mediana: 3.0
}%
\end{example}

\begin{example}
\exmp{
6
10 20 40 15 30 99
}{
Media: 35.7
Mediana: 25.0
}%
\end{example}

\begin{example}
\exmp{
8
1 1 2 1 1 1 1 1
}{
Media: 1.1
Mediana: 1.0
}%
\end{example}

%Comentários sobre o exemplo, se necessário.

\vspace{12pt}
Dica:

\lstset{language=C++,
                basicstyle=\ttfamily,
                keywordstyle=\color{Mulberry}\ttfamily,
                stringstyle=\color{Maroon}\ttfamily,
                commentstyle=\color{YellowGreen}\ttfamily,
                morecomment=[l][\color{magenta}]{\#},
                basewidth = {.48em}
}
\begin{lstlisting}
//Metodo bolha para ordenar um vetor
int Bolha (int a[], int n) {
    //Repita n-1 vezes
        //Para i = 0 ... n-2
            //Se a[i] > a[i+1]
                //Trocar a[i] com a[i+1]
}
\end{lstlisting}


\end{problem}

\end{document}
